\chapter{Hadamard Matrices}\label{conjectures:ch:hadamard}

\begin{center}\itshape
Thoughts on inverse orthogonal matrices, simultaneous sign-successions, and tessellated pavements in two or more colours, with applications to
Newton's rule, ornamental tile-work, and the theory of numbers.
\upshape(James Joseph Sylvester, the title of his paper of 1867)
\end{center}

Fill an $n\times n$ square with signs $+1$ and $-1$. How large can its determinant be? For $n=2$ the answer is $2$, from
$\bigl(\begin{smallmatrix}1&1\\1&-1\end{smallmatrix}\bigr)$; for $n=3$ it is $4$ (Exercise~\ref{conjectures:ex:det3}). In 1893 Jacques Hadamard proved that the
determinant of such a matrix is at most $n^{n/2}$, and that the bound is reached only when the rows are pairwise orthogonal. That can happen only when $n$
is $1$, $2$ or a multiple of $4$. Is it enough to be a multiple of $4$?

Nobody knows. Sylvester had already found, in his paper with the wonderful title above, examples of every order $2^k$, by a doubling trick. In 1933
Raymond Paley, a twenty-six-year-old Cambridge analyst working with Norbert Wiener at MIT, used the squares of finite fields to construct examples of
orders $q+1$ and $2(q+1)$ for prime powers $q$. His paper on ``orthogonal matrices'' appeared in 1933; on 7 April of that year he was killed by an
avalanche while skiing alone near Banff. John Williamson added a construction based on quaternions in 1944. With these tools the smallest order
left open was $92$, until Leonard Baumert, Solomon Golomb and Marshall Hall found a matrix of order $92$ in 1962, with Williamson's construction and
a computer at the Jet Propulsion Laboratory. The open orders kept falling, one by one, to cleverer constructions and longer searches: $428$ fell in
2004, to Hadi Kharaghani and Behruz Tayfeh-Rezaie, after twelve hours on a cluster of sixteen PCs. After that the smallest open order was $668$.

The \emph{Hadamard conjecture} says that a Hadamard matrix exists for every multiple of $4$. This chapter proves Hadamard's bound and the four classical
constructions, explains why a Hadamard matrix is the same thing as a certain block design and an optimal error-correcting code, and shows why the
problem is so stubborn: every known construction needs an algebraic miracle, and the miracles have gaps.

\begin{think}
Two vectors of $\pm1$ entries of length $n$ are orthogonal exactly when they agree in $\frac n2$ places. Try to write down four pairwise orthogonal
$\pm1$ vectors of length $4$, then three of length $6$. Where does the second attempt go wrong?
\end{think}

\section{Hadamard's bound}

We use Hadamard's inequality: for a real $n\times n$ matrix with rows $r_1,\dots,r_n$, $|\det A|\le\|r_1\|\cdots\|r_n\|$, with equality if and only if
some row is $0$ or the rows are pairwise orthogonal (Volume~8, Theorem~14.5). Geometrically: the volume of a parallelepiped is at most the product of its
edge lengths.

\begin{theorem}[Hadamard, 1893]\label{conjectures:thm:hadbound}
If $A$ is an $n\times n$ matrix with entries in $[-1,1]$, then $|\det A|\le n^{n/2}$. If the entries are $\pm1$, equality holds if and only if
$AA^{\mathsf T}=nI_n$.
\end{theorem}
\begin{proof}
Each row has $\|r_i\|\le\sqrt n$, so Hadamard's inequality gives $|\det A|\le n^{n/2}$. For a $\pm1$ matrix, $\|r_i\|=\sqrt n$ exactly and no row is $0$,
so equality holds if and only if the rows are pairwise orthogonal, that is $AA^{\mathsf T}=nI_n$.
\end{proof}

A \emph{Hadamard matrix} of order $n$ is an $n\times n$ matrix $H$ with entries $\pm1$ and $HH^{\mathsf T}=nI_n$. Then $H^{\mathsf T}=nH^{-1}$, so also
$H^{\mathsf T}H=nI_n$: the columns are orthogonal too. Permuting rows or columns, or multiplying some of them by $-1$, keeps the property; matrices related
in this way are called \emph{equivalent}. Every Hadamard matrix is equivalent to a \emph{normalised} one, whose first row and first column consist of $1$s.

\begin{proposition}\label{conjectures:prop:div4}
If a Hadamard matrix of order $n$ exists, then $n=1$, $n=2$ or $4\mid n$.
\end{proposition}
\begin{proof}
Let $n\ge3$ and normalise $H$, so that its first row consists of $1$s. Let $a,b,c,d$ be the numbers of columns whose entries in rows $2$ and $3$ are
$(+,+)$, $(+,-)$, $(-,+)$, $(-,-)$; then $a+b+c+d=n$. Orthogonality of row $1$ with rows $2$ and $3$ gives $a+b=c+d$ and $a+c=b+d$, and orthogonality
of rows $2$ and $3$ gives $a+d=b+c$. Subtracting the second equation from the first gives $b-c=c-b$, so $b=c$, and then $a=d$; the third equation
becomes $2a=2b$. So $a=b=c=d=\frac n4$, and $4\mid n$.
\end{proof}

The same argument appears as an exercise in Volume~3 (Exercise~K.4 there), and again in the problem book P1. The
size of the bound is worth a look: a random $\pm1$ matrix has $|\det A|$ of order roughly $\sqrt{n!}$, which is smaller than $n^{n/2}$ by a factor
of about $e^{n/2}$, so a Hadamard matrix is exponentially special.

\section{Sylvester's doubling}

The \emph{Kronecker product} of an $m\times m$ matrix $A$ and an $n\times n$ matrix $B$ is the $mn\times mn$ block matrix $A\otimes B=(a_{ij}B)$. It
satisfies $(A\otimes B)(C\otimes D)=(AC)\otimes(BD)$ and $(A\otimes B)^{\mathsf T}=A^{\mathsf T}\otimes B^{\mathsf T}$, by multiplying blocks.

\begin{proposition}[Sylvester, 1867]\label{conjectures:prop:kron}
If $H$ and $K$ are Hadamard matrices of orders $m$ and $n$, then $H\otimes K$ is a Hadamard matrix of order $mn$. In particular the matrices
$S_1=(1)$, $S_{2^{k+1}}=\begin{pmatrix}S_{2^k}&S_{2^k}\\S_{2^k}&-S_{2^k}\end{pmatrix}=\begin{pmatrix}1&1\\1&-1\end{pmatrix}\otimes S_{2^k}$ are
Hadamard matrices of order $2^k$.
\end{proposition}
\begin{proof}
$(H\otimes K)(H\otimes K)^{\mathsf T}=(HH^{\mathsf T})\otimes(KK^{\mathsf T})=(mI_m)\otimes(nI_n)=mnI_{mn}$, and the entries are products of signs.
\end{proof}

The rows of $S_{2^k}$ are the \emph{Walsh functions}: the entry in row $i$ and column $j$ is $(-1)^{i\cdot j}$, where $i\cdot j$ is the dot product of
the binary expansions (Exercise~\ref{conjectures:ex:walsh}). They are the characters of the group $(\mathbb Z/2)^k$, and multiplying by $S_{2^k}$ is the
Walsh--Hadamard transform of signal processing.

\section{Paley's constructions}

Let $q$ be an odd prime power and $\mathbb F_q$ the field with $q$ elements (Volume~3, Theorem~9.6). Its multiplicative group is cyclic (Volume~3,
Theorem~9.4); let $g$ be a generator. The \emph{quadratic character} $\chi:\mathbb F_q\to\{0,\pm1\}$ is $\chi(0)=0$, $\chi(x)=1$ if $x$ is a nonzero
square and $\chi(x)=-1$ otherwise.

\begin{lemma}\label{conjectures:lem:chi}
\begin{enumerate}
\item $\chi(g^j)=(-1)^j$, so $\chi$ is multiplicative and there are $\frac{q-1}2$ nonzero squares.
\item $\chi(-1)=(-1)^{(q-1)/2}$: $-1$ is a square if and only if $q\equiv1\pmod4$.
\item $\sum_{x\in\mathbb F_q}\chi(x)=0$.
\item For $c\neq0$, $\sum_{x\in\mathbb F_q}\chi(x)\chi(x+c)=-1$.
\end{enumerate}
\end{lemma}
\begin{proof}
(a) The squares of $g^i$ are the $g^{2i}$, so the nonzero squares are the even powers of $g$, which are $\frac{q-1}2$ elements since $g$ has even order
$q-1$. Hence $\chi(g^j)=(-1)^j$, and $\chi(g^ig^j)=(-1)^{i+j}=\chi(g^i)\chi(g^j)$; with $0$ the multiplicativity is clear. (b) $-1$ is the element of
order $2$, that is $g^{(q-1)/2}$. (c) There are as many $+1$ as $-1$. (d) The term $x=0$ is $0$. For $x\neq0$,
$\chi(x)\chi(x+c)=\chi(x)^2\chi(1+cx^{-1})=\chi(1+cx^{-1})$. As $x$ runs over $\mathbb F_q^*$, $1+cx^{-1}$ runs over $\mathbb F_q\setminus\{1\}$, so the sum is
$\sum_y\chi(y)-\chi(1)=0-1$.
\end{proof}

Index rows and columns by the elements of $\mathbb F_q$, in any order, and let $Q$ be the \emph{Jacobsthal matrix} $Q_{xy}=\chi(x-y)$. Write $J$ for an
all-ones matrix and $\mathbf j$ for an all-ones column, with sizes given by the context.

\begin{lemma}\label{conjectures:lem:jacobsthal}
$Q\mathbf j=0$, $\mathbf j^{\mathsf T}Q=0$, $QQ^{\mathsf T}=qI-J$, and $Q^{\mathsf T}=\chi(-1)Q$: the matrix $Q$ is symmetric if $q\equiv1\pmod4$ and
skew-symmetric if $q\equiv3\pmod4$.
\end{lemma}
\begin{proof}
The row and column sums are $\sum_y\chi(x-y)=0$ and $\sum_x\chi(x-y)=0$, by Lemma~\ref{conjectures:lem:chi}(c). The entry $(x,z)$ of $QQ^{\mathsf T}$ is
$\sum_y\chi(x-y)\chi(z-y)$; put $u=x-y$: it is $\sum_u\chi(u)\chi(u+z-x)$, which is $q-1$ if $z=x$ and $-1$ otherwise, by (a) and (d). Finally
$Q_{yx}=\chi(y-x)=\chi(-1)\chi(x-y)$.
\end{proof}

\begin{theorem}[Paley I, 1933]\label{conjectures:thm:paley1}
Let $q\equiv3\pmod4$ be a prime power and
\[
S=\begin{pmatrix}0&\mathbf j^{\mathsf T}\\-\mathbf j&Q\end{pmatrix}.
\]
Then $H=I+S$ is a Hadamard matrix of order $q+1$.
\end{theorem}
\begin{proof}
The entries of $S$ are $0$ on the diagonal and $\pm1$ elsewhere, because $\chi(x-y)=\pm1$ for $x\neq y$; so $H$ is a $\pm1$ matrix. Since $Q$ is
skew-symmetric, so is $S$. By Lemma~\ref{conjectures:lem:jacobsthal},
\[
SS^{\mathsf T}=\begin{pmatrix}0&\mathbf j^{\mathsf T}\\-\mathbf j&Q\end{pmatrix}\begin{pmatrix}0&-\mathbf j^{\mathsf T}\\\mathbf j&Q^{\mathsf T}\end{pmatrix}
=\begin{pmatrix}\mathbf j^{\mathsf T}\mathbf j&\mathbf j^{\mathsf T}Q^{\mathsf T}\\Q\mathbf j&\mathbf j\mathbf j^{\mathsf T}+QQ^{\mathsf T}\end{pmatrix}
=\begin{pmatrix}q&0\\0&J+qI-J\end{pmatrix}=qI.
\]
So $HH^{\mathsf T}=(I+S)(I-S)=I-S^2=I+SS^{\mathsf T}=(q+1)I$.
\end{proof}

\begin{theorem}[Paley II, 1933]\label{conjectures:thm:paley2}
Let $q\equiv1\pmod4$ be a prime power and
\[
S=\begin{pmatrix}0&\mathbf j^{\mathsf T}\\\mathbf j&Q\end{pmatrix},\qquad H=\begin{pmatrix}S+I&S-I\\S-I&-S-I\end{pmatrix}.
\]
Then $H$ is a Hadamard matrix of order $2(q+1)$.
\end{theorem}
\begin{proof}
Now $Q$ and $S$ are symmetric, and the computation of the previous proof, with $+\mathbf j$ in place of $-\mathbf j$, gives $S^2=SS^{\mathsf T}=qI$. The
entries of $S\pm I$ are $\pm1$, and $H$ is symmetric, so
\[
\begin{aligned}HH^{\mathsf T}=H^2&=\begin{pmatrix}(S+I)^2+(S-I)^2&(S+I)(S-I)-(S-I)(S+I)\\(S-I)(S+I)-(S+I)(S-I)&(S-I)^2+(S+I)^2\end{pmatrix}\\
&=\begin{pmatrix}2(q+1)I&0\\0&2(q+1)I\end{pmatrix},\end{aligned}
\]
since $S$ commutes with $I$ and $(S+I)^2+(S-I)^2=2S^2+2I$.
\end{proof}

The matrix $S$ of Paley~I is a \emph{conference matrix}: zero diagonal, $\pm1$ elsewhere, $SS^{\mathsf T}=(n-1)I$. Such matrices were studied for
telephone conference networks, hence the name. Figure~\ref{conjectures:fig:paley} shows two Paley matrices. With Kronecker products, Paley's two constructions give
Hadamard matrices of all orders $4,8,\dots,88$; the first order they miss is $92$, and up to $1000$ they miss $55$ multiples of $4$ (Python
companion, P7.2).

\begin{figure}[ht]
\centering
\fig{conjectures-ch07-paley}
\caption{Two Paley matrices, with $-1$ in black. Left: Paley~I for $q=43$, of order $44$; below the first row and column, the entry in row $x$ and
column $y$ is $\chi(x-y)$, so the pattern is constant along diagonals and the matrix is skew apart from its diagonal. Right: Paley~II for $q=13$, of
order $28$, made of four blocks $S\pm I$. Both were checked to satisfy $HH^{\mathsf T}=nI$.}
\label{conjectures:fig:paley}
\end{figure}

\section{Williamson's construction}

Euler's four-squares identity says that a product of two sums of four squares is a sum of four squares; behind it is the multiplication of
quaternions, $|xy|=|x||y|$. Williamson turned the identity into matrices.

\begin{theorem}[Williamson, 1944]\label{conjectures:thm:williamson}
Let $A,B,C,D$ be symmetric $n\times n$ matrices with entries $\pm1$ that commute pairwise and satisfy $A^2+B^2+C^2+D^2=4nI$. Then
\[
W=\begin{pmatrix}A&B&C&D\\-B&A&-D&C\\-C&D&A&-B\\-D&-C&B&A\end{pmatrix}
\]
is a Hadamard matrix of order $4n$.
\end{theorem}
\begin{proof}
Write the block rows as $R_1,\dots,R_4$ and compute the blocks $R_iR_k^{\mathsf T}=\sum_mX_mY_m^{\mathsf T}$, using $Y^{\mathsf T}=Y$ for the four
matrices. Each diagonal block is $A^2+B^2+C^2+D^2=4nI$. Above the diagonal,
\begin{align*}
R_1R_2^{\mathsf T}&=-AB+BA-CD+DC,\\ R_1R_3^{\mathsf T}&=-AC+BD+CA-DB,\\
R_1R_4^{\mathsf T}&=-AD-BC+CB+DA,\\ R_2R_3^{\mathsf T}&=BC+AD-DA-CB,\\
R_2R_4^{\mathsf T}&=BD-AC-DB+CA,\\ R_3R_4^{\mathsf T}&=CD-DC+AB-BA.
\end{align*}
Each is a sum of two commutators $XY-YX$, which vanish. The blocks below the diagonal are the transposes of these. So $WW^{\mathsf T}=4nI$.
\end{proof}

The four blocks are the coefficients of a quaternion $a+bi+cj+dk$, and $W$ is the matrix of left multiplication by it, with $1,i,j,k$ replaced by
matrices: the rows of that matrix are orthogonal because of the norm identity, which is Euler's identity. In practice one takes $A,B,C,D$
\emph{circulant}, determined by their first rows; circulant matrices commute automatically, and the condition becomes a statement about the
autocorrelations of four short sequences. For $n=3$, $A=J$ and $B=C=D=2I-J$ work (Exercise~\ref{conjectures:ex:williamson3}), which gives order $12$. Baumert,
Golomb and Hall found Williamson matrices of order $23$, hence a Hadamard matrix of order $92$. They do not exist for every $n$: Djokovi\'c showed by an
exhaustive search in 1993 that there are none for $n=35$, although a Hadamard matrix of order $140$ comes at once from Paley~I with $q=139$.

\section{Designs and codes}

A Hadamard matrix hides two other objects. A \emph{symmetric $2$-$(v,k,\lambda)$ design} is a set of $v$ points with $v$ subsets called blocks, each of
$k$ points, such that every pair of distinct points lies in exactly $\lambda$ blocks. Counting the pairs $(y,B)$ with $x,y\in B$, $y\neq x$, for a fixed
point $x$ shows that $x$ lies in $r=\frac{\lambda(v-1)}{k-1}$ blocks.

\begin{theorem}\label{conjectures:thm:haddesign}
Let $t\ge2$. Hadamard matrices of order $4t$ correspond to symmetric $2$-$(4t-1,2t-1,t-1)$ designs: from a normalised Hadamard matrix, take as points
the rows $2,\dots,4t$ and, for each column $j\ge2$, the block of rows $i\ge2$ with $H_{ij}=1$; conversely, from the $0/1$ incidence matrix $N$ of such a
design, the matrix
\[
H=\begin{pmatrix}1&\mathbf j^{\mathsf T}\\\mathbf j&2N-J\end{pmatrix}
\]
is a Hadamard matrix.
\end{theorem}
\begin{proof}
From $H$ to the design: column $j\ge2$ is orthogonal to the first column, so it has $2t$ entries $+1$, one of them in row $1$; the block has $2t-1$
points. For two rows $i\neq k$, both $\ge2$, count the columns by the signs in rows $i$ and $k$: as in Proposition~\ref{conjectures:prop:div4}, all four patterns
occur $t$ times, so $(+,+)$ occurs in $t$ columns, one of which is column $1$. The pair lies in $t-1$ blocks.

Conversely let $N$ be the $v\times v$ incidence matrix ($N_{xB}=1$ if $x\in B$), with $v=4t-1$, $k=2t-1$, $\lambda=t-1$. Then $NN^{\mathsf T}$ has
$r$ on the diagonal and $\lambda$ elsewhere, and $r=\frac{(t-1)(4t-2)}{2t-2}=2t-1$; so $NN^{\mathsf T}=tI+(t-1)J$, and $NJ=JN^{\mathsf T}=(2t-1)J$
because every row and every column of $N$ has $2t-1$ ones. With $M=2N-J$,
\[
MM^{\mathsf T}=4NN^{\mathsf T}-2NJ-2JN^{\mathsf T}+J^2=4tI+\bigl(4(t-1)-4(2t-1)+(4t-1)\bigr)J=4tI-J,
\]
and $M\mathbf j=(2(2t-1)-(4t-1))\mathbf j=-\mathbf j$. So
\[
HH^{\mathsf T}=\begin{pmatrix}4t&\mathbf j^{\mathsf T}+\mathbf j^{\mathsf T}M^{\mathsf T}\\\mathbf j+M\mathbf j&J+MM^{\mathsf T}\end{pmatrix}=4tI.
\]
\end{proof}

For $t=2$ the design is a $2$-$(7,3,1)$ design, the \emph{Fano plane} (Figure~\ref{conjectures:fig:fano}); for $t=3$ it is a $2$-$(11,5,2)$ design, a ``biplane''.

\begin{figure}[ht]
\centering
\fig{conjectures-ch07-fano}
\caption{The Fano plane obtained from the Sylvester matrix $S_8$ by Theorem~\ref{conjectures:thm:haddesign}. The points are the rows $1,\dots,7$ of $S_8$ (counted
from $0$), the lines are its columns $1,\dots,7$, and a point is on a line when the entry is $+1$. The seven lines are the six segments and the
circle; any two points lie on exactly one line.}
\label{conjectures:fig:fano}
\end{figure}

For codes, replace $+1$ by $0$ and $-1$ by $1$. Two $\pm1$ vectors $u,v$ of length $n$ differ in $d(u,v)=\frac{n-\langle u,v\rangle}2$ places.

\begin{proposition}\label{conjectures:prop:hadcode}
The $2n$ rows of $H$ and $-H$, for a Hadamard matrix $H$ of order $n$, are binary words of length $n$ any two of which differ in at least $\frac n2$
places. No binary code of length $n$ with this property has more than $2n$ words.
\end{proposition}
\begin{proof}
Two distinct rows of $H$, or of $-H$, are orthogonal, so they differ in $\frac n2$ places; a row of $H$ and a row of $-H$ differ in $\frac n2$ places, or in
$n$ if one is the negative of the other. For the bound, a code with distance at least $\frac n2$ gives $\pm1$ vectors with
$\langle u,v\rangle=n-2d(u,v)\le0$ for $u\neq v$, and the next lemma applies.
\end{proof}

\begin{lemma}\label{conjectures:lem:obtuse}
At most $2n$ nonzero vectors of $\mathbb R^n$ can have pairwise nonpositive inner products.
\end{lemma}
\begin{proof}
Induction on $n$. For $n=1$, two nonzero numbers with nonpositive product have opposite signs, so there are at most two. Let $n\ge2$ and let
$v_1,\dots,v_M$ be such vectors. Project the others onto the hyperplane $v_1^\perp$: $w_i=v_i-\frac{\langle v_i,v_1\rangle}{\|v_1\|^2}v_1$. Then for
$i\neq k$, $\langle w_i,w_k\rangle=\langle v_i,v_k\rangle-\frac{\langle v_i,v_1\rangle\langle v_k,v_1\rangle}{\|v_1\|^2}\le0$, since the subtracted
product of two nonpositive numbers is nonnegative. A $w_i$ is $0$ only if $v_i$ is a negative multiple of $v_1$, and two such $v_i$ would have a positive
inner product; so at least $M-2$ of the $w_i$ are nonzero vectors with pairwise nonpositive inner products in an $(n-1)$-dimensional space. By
induction $M-2\le2(n-1)$.
\end{proof}

The code of Proposition~\ref{conjectures:prop:hadcode} built from $S_{2^k}$ is the first-order Reed--Muller code. Its decoder is a single Walsh--Hadamard
transform: correlate the received word with every row and pick the largest correlation.

\begin{remark}[other orders, other counts]\label{conjectures:rem:barba}
When $4\nmid n$ the bound $n^{n/2}$ is not attained, and better bounds are known: $|\det A|\le\sqrt{2n-1}\,(n-1)^{(n-1)/2}$ for odd $n$ (Guido Barba,
1933) and $|\det A|\le2(n-1)(n-2)^{(n-2)/2}$ for $n\equiv2\pmod4$ (Hans Ehlich, 1964). There are many Hadamard matrices: Kharaghani and Tayfeh-Rezaie
showed in 2010 that there are exactly $13\,710\,027$ inequivalent ones of order $32$. And asymptotically the conjecture is ``almost'' true: Jennifer
Seberry proved in 1976 that for every odd $q$ there is a Hadamard matrix of order $2^tq$ for all $t$ larger than about $2\log_2q$.
\end{remark}

\section{Why it is hard}

\paragraph{Every construction is a miracle.} Sylvester uses the group $(\mathbb Z/2)^k$, Paley a finite field, Williamson quaternions and sequences with
vanishing autocorrelations. Each needs a structure that exists only for special $n$, and the products and combinations of known matrices cover a set
of orders that grows, but has always had gaps. No construction is known that takes $n$ as input and works for all multiples of $4$.

\paragraph{Counting does not help.} There are $2^{n^2}$ sign matrices, and the Hadamard ones are an exponentially small fraction: by
Theorem~\ref{conjectures:thm:hadbound} they have the largest possible determinant, while a typical one has a determinant smaller by a factor of about $e^{n/2}$.
A probabilistic argument would have to find a needle in this haystack, and nobody knows how to set one up. At the same time there are a great many
needles (Remark~\ref{conjectures:rem:barba}), which is why the conjecture is believed.

\paragraph{Searches grow fast.} The open orders are attacked by searching for short sequences with special autocorrelations, which then feed a
construction. The search space grows exponentially with the length of the sequences. The order $428$ needed four sequences of lengths $36,36,36,35$,
found after twelve hours on sixteen machines; each new order is a new search.

\begin{remark}[the status of $668$]\label{conjectures:rem:668}
From 2004 the smallest order without a known Hadamard matrix was $668=4\cdot167$, followed by $716$ and $892$ (the order $764$ was settled by Dragomir
Djokovi\'c in 2008). In 2026 a group of researchers announced Hadamard matrices for all the open orders up to $2000$, including $668$; at the time of
writing, the matrices had been posted but no paper had appeared. Whatever the outcome, the conjecture itself, for all multiples of $4$, remains open.
\end{remark}

\section{Reading a paper: the order 428}

Hadi Kharaghani and Behruz Tayfeh-Rezaie, \emph{A Hadamard matrix of order 428} (Journal of Combinatorial Designs, 2005), is six pages long. It is a
chain of three classical constructions, and the paper is a good place to see how the field works. Four $\pm1$ sequences of lengths $36,36,36,35$ with
a condition on their aperiodic autocorrelations (\emph{Turyn-type} sequences) are found by computer; a theorem of Turyn turns them into \emph{base
sequences} of lengths $71,71,36,36$; these give \emph{T-sequences} of length $107$; and four circulant matrices built from the T-sequences, placed in
an array of the Williamson type with a reversal matrix, give a Hadamard matrix of order $4\cdot107=428$. Read the paper for the chain, and do not try to
reproduce the search: it needed $1\,911\,620$ partial solutions before the full one. The idea worth stealing is the reduction of a matrix of size $428^2$
to four sequences of length about $36$, a search space small enough for a computer.

\section{Open doors}

\begin{itemize}
\item List the multiples of $4$ up to $2000$ not covered by Sylvester's and Paley's constructions and their Kronecker products. Which of them are
covered by Williamson matrices of the orders you can find by computer?
\item Search for symmetric circulant Williamson matrices of every odd order $n\le25$. For which $n$ do they exist? Compare with the literature.
\item Compute the largest determinant of a $\pm1$ matrix of order $n$ for $n\le7$ by a clever search, and compare with Barba's and Ehlich's bounds.
\item Build the Hadamard code of length $32$ and measure, by simulation, how many random bit errors its Walsh--Hadamard decoder corrects.
\end{itemize}

\section{Exercises}

\begin{exercise}[$\star$]\label{conjectures:ex:order4}
Show that any two Hadamard matrices of order $4$ are equivalent.
\end{exercise}

\begin{exercise}[$\star\star$]\label{conjectures:ex:det3}
Show that the determinant of an $n\times n$ matrix with entries $\pm1$ is divisible by $2^{n-1}$. Deduce that the largest determinant of a $3\times3$
sign matrix is $4$, and find a matrix that attains it.
\end{exercise}

\begin{exercise}[$\star$]\label{conjectures:ex:walsh}
Number the rows and columns of $S_{2^k}$ from $0$ to $2^k-1$ and show that its entry $(i,j)$ is $(-1)^{i\cdot j}$, where $i\cdot j$ is the number of
positions where the binary expansions of $i$ and $j$ both have a $1$. Show that $S_{2^k}x$ can be computed with $k2^k$ additions and subtractions.
\end{exercise}

\begin{exercise}[$\star$]\label{conjectures:ex:paley5}
Write down the Jacobsthal matrix for $q=5$, check that it is symmetric with $Q^2=5I-J$, and write the conference matrix $S$ of order $6$ of
Theorem~\ref{conjectures:thm:paley2}. What is the order of the resulting Hadamard matrix?
\end{exercise}

\begin{exercise}[$\star\star$]\label{conjectures:ex:charsum}
Let $q$ be an odd prime power and $a\in\mathbb F_q^*$. By counting the solutions of $y^2=x^2-a$, show that $\sum_{x\in\mathbb F_q}\chi(x^2-a)=-1$.
\end{exercise}

\begin{exercise}[$\star$]\label{conjectures:ex:williamson3}
Check that $A=J_3$ and $B=C=D=2I_3-J_3$ satisfy the hypotheses of Theorem~\ref{conjectures:thm:williamson}, and write the resulting Hadamard matrix of order $12$.
\end{exercise}

\begin{exercise}[$\star\star$]\label{conjectures:ex:williamsonsums}
Let $A,B,C,D$ be Williamson matrices of odd order $n$ that are circulant, with row sums $a,b,c,d$. Show that $a^2+b^2+c^2+d^2=4n$ and that $a,b,c,d$ are
odd. For $n=5$, which row sums are possible?
\end{exercise}

\begin{exercise}[$\star$]\label{conjectures:ex:fano}
Check, from the matrix $S_8$, the seven lines of Figure~\ref{conjectures:fig:fano}, and that any two points lie on exactly one line.
\end{exercise}

\begin{exercise}[$\star\star$]\label{conjectures:ex:obtuseeq}
Show that if $2n$ nonzero vectors of $\mathbb R^n$ have pairwise nonpositive inner products, they are $\pm$ multiples of the vectors of an orthogonal
basis: they come in $n$ pairs $u_i$, $-c_iu_i$ ($c_i>0$) with $u_1,\dots,u_n$ pairwise orthogonal.
\end{exercise}

\begin{exercise}[$\star$]\label{conjectures:ex:biplane}
Find the parameters of the design of Theorem~\ref{conjectures:thm:haddesign} for $t=3$ and check the relation $r(k-1)=\lambda(v-1)$. How many blocks contain two
given points, and how many points do two blocks share? (Use $N^{\mathsf T}N$.)
\end{exercise}

\section{Hints}

\begin{hint}{conjectures:ex:order4}
Normalise, then use Proposition~\ref{conjectures:prop:div4}: each of rows $2$, $3$, $4$ has two $+1$ and two $-1$, and they are orthogonal.
\end{hint}
\begin{hint}{conjectures:ex:det3}
Subtract the first row from the others, then take a factor $2$ from each of the other $n-1$ rows. For $n=3$, $4\mid\det$ and $|\det|\le3^{3/2}$.
\end{hint}
\begin{hint}{conjectures:ex:walsh}
Induction on $k$: the block structure of $S_{2^{k+1}}$ is decided by the top bits of $i$ and $j$. For the speed, $S_{2^{k+1}}x$ needs two products by
$S_{2^k}$ and $2^{k+1}$ additions.
\end{hint}
\begin{hint}{conjectures:ex:paley5}
The nonzero squares modulo $5$ are $1$ and $4$.
\end{hint}
\begin{hint}{conjectures:ex:charsum}
$y^2=x^2-a$ means $(x-y)(x+y)=a$. And the number of $y$ with $y^2=z$ is $1+\chi(z)$.
\end{hint}
\begin{hint}{conjectures:ex:williamson3}
$(2I-J)^2=4I-4J+J^2$ and $J_3^2=3J_3$.
\end{hint}
\begin{hint}{conjectures:ex:williamsonsums}
Multiply $A^2+B^2+C^2+D^2=4nI$ by $\mathbf j$. A row of odd length with entries $\pm1$ has an odd sum.
\end{hint}
\begin{hint}{conjectures:ex:fano}
Row $i$ of $S_8$ has $+1$ in column $j$ when $i\cdot j$ is even.
\end{hint}
\begin{hint}{conjectures:ex:obtuseeq}
Follow the proof of Lemma~\ref{conjectures:lem:obtuse}: equality forces one $v_i$ to be a negative multiple of $v_1$. What does that say about $\langle v_k,v_1\rangle$
for the others?
\end{hint}
\begin{hint}{conjectures:ex:biplane}
$t=3$. For the last question, $N^{\mathsf T}$ is the incidence matrix of the dual structure; compute $N^{\mathsf T}N$ from $NN^{\mathsf T}$ and the
invertibility of $N$.
\end{hint}

\section{Solutions}

\begin{solution}{conjectures:ex:order4}
\idea Normalisation leaves almost no freedom.
\writeup Normalise: first row and column all $1$. Each of rows $2,3,4$ is orthogonal to row $1$, so it has two $+1$ and two $-1$, with a $+1$ in column
$1$: it is $(1,1,-1,-1)$, $(1,-1,1,-1)$ or $(1,-1,-1,1)$. Rows $2,3,4$ are distinct, because two equal rows are not orthogonal, so they are these three
vectors in some order; permuting rows $2,3,4$ brings any normalised matrix to the same one. Every Hadamard matrix of order $4$ is therefore equivalent to
it.
\end{solution}

\begin{solution}{conjectures:ex:det3}
\idea Row operations make the other rows even.
\writeup Subtracting row $1$ from each other row does not change the determinant, and the new rows $2,\dots,n$ have entries in $\{0,\pm2\}$. Taking
a factor $2$ out of each of these $n-1$ rows leaves an integer matrix, so $2^{n-1}\mid\det A$. For $n=3$: $4\mid\det A$ and
$|\det A|\le3^{3/2}=5.196$, so $|\det A|\le4$. The matrix with rows $(1,1,1)$, $(1,-1,1)$, $(1,1,-1)$ has determinant
$1\cdot(1-1)-1\cdot(-1-1)+1\cdot(1+1)=4$.
\end{solution}

\begin{solution}{conjectures:ex:walsh}
\idea Doubling adds one bit.
\writeup For $k=0$ the entry is $1=(-1)^0$. Assume the formula for $S_{2^k}$. In $S_{2^{k+1}}$ write $i=i_k2^k+i'$, $j=j_k2^k+j'$ with $i',j'<2^k$. The
block in position $(i_k,j_k)$ is $S_{2^k}$, except the block $(1,1)$, which is $-S_{2^k}$; so the entry is $(-1)^{i_kj_k}(-1)^{i'\cdot j'}=(-1)^{i\cdot j}$.
For the speed, write $x=(x',x'')$ in halves: $S_{2^{k+1}}x=(S_{2^k}x'+S_{2^k}x'',S_{2^k}x'-S_{2^k}x'')$, so the cost $c_{k+1}$ satisfies
$c_{k+1}=2c_k+2^{k+1}$ with $c_0=0$, which gives $c_k=k2^k$.
\end{solution}

\begin{solution}{conjectures:ex:paley5}
\idea $\chi$ modulo $5$ is $\chi(1)=\chi(4)=1$, $\chi(2)=\chi(3)=-1$.
\writeup Number the rows and columns $0,\dots,4$. Row $0$ lists the values $\chi(-y)$, and each next row is the previous one shifted right by one:
\[
Q=\begin{pmatrix}0&1&-1&-1&1\\1&0&1&-1&-1\\-1&1&0&1&-1\\-1&-1&1&0&1\\1&-1&-1&1&0\end{pmatrix}.
\]
It is symmetric because $\chi(-1)=1$. Each row has four nonzero entries, so the diagonal of $Q^2$ is $4$; two distinct rows, for instance the first two, have inner product
$0+0-1+1-1=-1$, and the others in the same way, so $Q^2=4I-(J-I)=5I-J$. The conference matrix is the matrix $S$ of order $6$ with first row
$(0,1,1,1,1,1)$, first column the same, and $Q$ in the remaining block; $S^2=5I$, and Theorem~\ref{conjectures:thm:paley2} gives a Hadamard matrix of order $12$.
\end{solution}

\begin{solution}{conjectures:ex:charsum}
\idea Count the points of a hyperbola in two ways.
\writeup Put $u=x-y$, $w=x+y$; since $2$ is invertible, $(x,y)\mapsto(u,w)$ is a bijection of $\mathbb F_q^2$, and $y^2=x^2-a$ becomes $uw=a$, which has
$q-1$ solutions ($u\neq0$ arbitrary, $w=a/u$). On the other hand, for each $x$ the number of $y$ with $y^2=x^2-a$ is $1+\chi(x^2-a)$: two if $x^2-a$ is a
nonzero square, one if it is $0$, none otherwise. So $q-1=\sum_x(1+\chi(x^2-a))=q+\sum_x\chi(x^2-a)$.
\end{solution}

\begin{solution}{conjectures:ex:williamson3}
\idea Everything is a polynomial in $J$.
\writeup All four matrices are symmetric and are polynomials in $J_3$, so they commute. $A^2=J^2=3J$ and $B^2=(2I-J)^2=4I-4J+3J=4I-J$, so
$A^2+3B^2=3J+12I-3J=12I=4\cdot3\,I$. The Hadamard matrix is the $12\times12$ array of Theorem~\ref{conjectures:thm:williamson}: its first block row is
$(J,\,2I-J,\,2I-J,\,2I-J)$, where $2I-J$ has $1$ on the diagonal and $-1$ elsewhere, and the other block rows follow the pattern of signs.
\end{solution}

\begin{solution}{conjectures:ex:williamsonsums}
\idea A circulant matrix has a constant row sum, and $X\mathbf j$ is that sum times $\mathbf j$.
\writeup $A\mathbf j=a\mathbf j$, so $A^2\mathbf j=a^2\mathbf j$, and the same for $B,C,D$. Applying $A^2+B^2+C^2+D^2=4nI$ to $\mathbf j$ gives
$a^2+b^2+c^2+d^2=4n$. A sum of $n$ terms $\pm1$, $n$ odd, is odd. For $n=5$: $20$ as a sum of four odd squares ($1$, $9$; $25>20$) is only $9+9+1+1$, so
up to order and signs the row sums are $3,3,1,1$.
\end{solution}

\begin{solution}{conjectures:ex:fano}
\idea Apply Theorem~\ref{conjectures:thm:haddesign} and Exercise~\ref{conjectures:ex:walsh}.
\writeup Point $i$ ($1\le i\le7$) is on line $j$ when $i\cdot j$ is even. Line $1$ (binary $001$) contains the $i$ with last bit $0$: $2,4,6$. Line $2$
($010$): $1,4,5$. Line $3$ ($011$): $i$ with the last two bits equal: $3,4,7$. Line $4$ ($100$): $1,2,3$. Line $5$ ($101$): $2,5,7$. Line $6$ ($110$):
$1,6,7$. Line $7$ ($111$): $i$ with an even number of $1$s: $3,5,6$. These are the six segments and the circle of the figure. Each of the $21$ pairs of
points appears in exactly one of the seven lines, which hold $7\cdot3=21$ pairs in total, so no pair is repeated.
\end{solution}

\begin{solution}{conjectures:ex:obtuseeq}
\idea Equality in the induction forces a negative partner for every vector, and orthogonality to everything else.
\writeup Induction on $n$; for $n=1$ the two vectors are $u$ and $-cu$. For $n\ge2$ and $M=2n$, the proof of Lemma~\ref{conjectures:lem:obtuse} gives
$M-2\le$ the number of nonzero $w_i$, and $M-2=2n-2$ is the maximum possible in dimension $n-1$; so exactly one $v_i$, say $v_2=-cv_1$, is a negative
multiple of $v_1$. For $k\ge3$, $\langle v_k,v_1\rangle\le0$ and $\langle v_k,v_2\rangle=-c\langle v_k,v_1\rangle\le0$, so $\langle v_k,v_1\rangle=0$:
all the other vectors lie in $v_1^\perp$, where there are $2n-2$ of them, and induction applies.
\end{solution}

\begin{solution}{conjectures:ex:biplane}
\idea Symmetric designs are self-dual.
\writeup $t=3$: $v=11$, $k=5$, $\lambda=2$, and $r=\frac{2\cdot10}{4}=5$, so $r(k-1)=20=\lambda(v-1)$. Two points lie in $2$ blocks. For two blocks:
$NN^{\mathsf T}=3I+2J$ is invertible, hence so is $N$, and $NJ=JN=5J$. Then $N^{\mathsf T}N=N^{-1}(NN^{\mathsf T})N=N^{-1}(3I+2J)N=3I+2N^{-1}JN=3I+2J$,
because $JN=5J$ and $N^{-1}J=\frac15J$. So two blocks share exactly $2$ points: the dual of a symmetric design is again one, with the same parameters.
\end{solution}

\section*{Sources and further reading}

\begin{itemize}
\item J.~J.~Sylvester, \emph{Thoughts on inverse orthogonal matrices, simultaneous sign-successions, and tessellated pavements in two or more colours,
with applications to Newton's rule, ornamental tile-work, and the theory of numbers}, Philosophical Magazine 34 (1867), 461--475.
\item J.~Hadamard, \emph{R\'esolution d'une question relative aux d\'eterminants}, Bulletin des Sciences Math\'ematiques 17 (1893), 240--246.
\item R.~E.~A.~C.~Paley, \emph{On orthogonal matrices}, Journal of Mathematics and Physics 12 (1933), 311--320. For Paley's life and death: the MacTutor
biography of Raymond Paley, School of Mathematics and Statistics, University of St Andrews.
\item J.~Williamson, \emph{Hadamard's determinant theorem and the sum of four squares}, Duke Mathematical Journal 11 (1944), 65--81.
\item L.~Baumert, S.~W.~Golomb and M.~Hall, \emph{Discovery of an Hadamard matrix of order 92}, Bulletin of the AMS 68 (1962), 237--238.
\item H.~Kharaghani and B.~Tayfeh-Rezaie, \emph{A Hadamard matrix of order 428}, Journal of Combinatorial Designs 13 (2005), 435--440.
H.~Kharaghani and B.~Tayfeh-Rezaie, \emph{On the classification of Hadamard matrices of order 32}, Journal of Combinatorial Designs 18 (2010), 328--336.
\item D.~Z.~Djokovi\'c, Discrete Mathematics 115 (1993), 267--271 (the non-existence of Williamson matrices for $n=35$).
\item D.~Z.~Djokovi\'c, \emph{Hadamard matrices of order 764 exist}, Combinatorica (2008); arXiv:\allowbreak math/0703312.
\item J.~Seberry Wallis, \emph{On the existence of Hadamard matrices}, Journal of Combinatorial Theory A 21 (1976); and the correction of its bound in
arXiv:0706.1720.
\item K.~J.~Horadam, \emph{Hadamard Matrices and Their Applications}, Princeton University Press, 2007. J.~Seberry and M.~Yamada, \emph{Hadamard
Matrices: Constructions using Number Theory and Linear Algebra}, Wiley, 2020.
\item On the 2026 announcement for the open orders up to $2000$: Epoch AI, \emph{FrontierMath open problems: Hadamard matrices}, web page (accessed
October 2026).
\end{itemize}
